% Preamble.
\documentclass{article}

% Packages.
\usepackage[utf8]{inputenc}     % Unicode support (Umlauts etc.)
\usepackage{amsmath, amssymb}   % Advanced math typesetting
\usepackage{stmaryrd, bbm}
\usepackage[french]{babel}      % Change hyphenation rules
\usepackage[T1]{fontenc}
\usepackage{hyperref}           % Add a link to your document
\usepackage{graphicx}           % Add pictures to your document
\usepackage{listings}           % Source code formatting and highlighting
\usepackage{csquotes}
\usepackage{interval}

% Macros.
%% Double interval.
\newcommand{\iinterval}[2]{\llbracket #1; #2 \rrbracket}
%% Finite set, with curly braces.
\newcommand{\setf}[1]{\{#1 \}}
%% Floor brackets.
\newcommand{\floor}[1]{\lfloor #1 \rfloor}
%% Ceiling brackets.
\newcommand{\ceil}[1]{\lceil #1 \rceil}
%% Pre-formatted "Exemple".
\newcommand{\Ex}{\textsc{Exemple}: }

% Data.
\title{Terminale - Notations mathématiques.}
\author{Hector GARLATTI}
\date{\today{}}

\begin{document}
\sloppy 
\maketitle

\section{Utilité.}
Ce document est une fiche, rédigée en \LaTeX, d'élève de fin de Terminale. 
Surtout destinée aux révisions pour les classes préparatoires, 
elle peut servir d'aperçu pour les élèves de Première.

\subsection{Introduction.}
Le langage mathématique permet une expression objective, non-ambiguë
des théorèmes et des propriétés. Je pense qu'il est déjà utile en début
de Terminale de se familiariser avec ses notions car elles permettent:
\begin{itemize}
    \item la précision et la rigueur;
    \item de commencer à s'immerger dans l'univers
    des "mathématiques sérieuses";
    \item exprimer des concepts plus avancés.
\end{itemize}

Cependant, ces notations ne sont pas de excuses pour ne pas rédiger
et sont d'abord des outils pour enrichir la rédaction, non la simplifier.

\subsection{Règles pour la démonstration.}
En fin de Terminale, certains règles précises doivent être maîtrisés. 
En voici quelques unes:
\begin{enumerate}
    \item Un discours mathématiques n'est pas un suite de symboles.
    \item Tout objet est \textit{déclaré} avant d'être utilisé. \\
        \Ex \enquote{\underline{Soit} $x$ un nombre réel positif.}
    \item Une bonne démonstration est structurée par des locutions 
    qui guident le lecteur, que nous pouvons mettre en évidence
    (souligné, italique). \\
        \Ex \enquote{
            \underline{Montrons que} ... 
            \underline{Ainsi} ...
        }
\end{enumerate}


\section{Notations, quantificateurs.}
Ici ne seront que listé et brièvement mis en contexte les notation.
Pour des explications plus approfondies, se référer au leçons.

\subsection{Ensembles.}

\paragraph{Ensembles de nombres usuels.}
\begin{itemize}
    \item Notons un ensemble: $E = \{a; b; c; ...\}$:
    \begin{itemize}
        \item son nombre d'éléments: $\text{Card}({E}) = |E|$;
        \item l'ordre ne compte pas; \\
        \Ex $\{1; 2\} = \{2; 1\}$
        \item les répétitions sont "effacées"; \\
        \Ex $\{1; 2; 2\} = \{1; 2\}$
    \end{itemize} 
    \item $\varnothing = \{\}$: l'ensemble vide.
    \item $\mathbb{N}$: les entiers naturels, \\
        \Ex $\{0; 1; 2; \dots\}$. 
    \item $\mathbb{Z}$: les entiers relatifs, \\ 
        \Ex $\{0; 1; -1; 2; -2; ...\}$.
    \item $\mathbb{Q}$: les nombres rationnels, toutes les fractions 
    $$
    \forall (p; q) \in \mathbb{Z}\times \mathbb{N}^{\star} \quad 
    \frac{p}{q}
    $$.
    \item $\mathbb{R}$: le continuum des réels, \\
        \Ex $\{\sqrt{2}; \pi; \sin(2); \dots\}$.
    \item $\mathbb{C}$: les complexes, \\
        \Ex $\{1 + 2i; 2e^{\frac{1}{4}i\pi}; \dots\}$.
    \item $\mathbb{R}^{n}$: les vecteurs (vecteurs, points, $\dots$) de $n$ dimensions. \\
        \Ex $(1; 2) \in \mathbb{R}^2$
    \item $E[X]$: les polynômes à valeurs dans $E$. \\
        \Ex $(x^2 + 3) \in \mathbb{R}[X]$ 
    \item $\mathcal{M}_{n, m}(E)$: les matrices de taille $(n, m)$ 
        à coefficients dans $E$ \\
        \Ex $\mathbbm{1}_3 \in \mathcal{M}_{3}(\mathbb{N})$
\end{itemize}

\paragraph{Appartenance.} 
Notations permettant de déclarer un object (\textit{cf.} \textbf{1.2.3})
de le décrire.

\begin{itemize}
    \item $\in$: un scalaire, un nombre, un vecteur fait parti d'un ensemble $E$. \\
        \Ex 
        \enquote{$x \in \mathbb{R}$} signifie 
        \enquote{\underline{Soit} $x$ un nombre réel.}
    \item $\subset$: un ensemble $A$ est compris dans, est un sous-ensemble de $B$. \\
        \Ex $A \subset B$ \\
        Ainsi, nous pouvons écrire pour nos ensembles usuels 
        (penser au diagramme des cercles): 
        \begin{equation}
            \mathbb{N} \subset \mathbb{Z} \subset \mathbb{Q} 
            \subset \mathbb{R} \subset \mathbb{C}
        \end{equation}
\end{itemize}

\paragraph{Altérations et operations.} 
Pour un ensemble $E$.
\begin{itemize}
    \item $E^{\star}$: excluant $0$. \\
        \Ex $\mathbb{N}^{\star}$, tout les entiers naturel 
        \textit{hormis} $0$.
    \item $E_{+}$: gardant uniquement les positifs. \\
        \Ex $\mathbb{R}_{+}$, tout les réels \textit{positifs}.
    \item $E_{-}$: gardant uniquement les négatifs. \\
        \Ex $\mathbb{R}^{\star}_{-}$, tout les réels négatifs sans $0$.
    
    \item $\setminus$: difference d'ensembles. \\
        \Ex $\mathbb{R} \setminus \{1\}$, tout les réels excluant $1$.
    \item $\cup$: union, \enquote{ou inclusif} logique. \\
        \Ex $\{1; 2\} \cup \{2; 3\} = \{1; 2; 3\}$
    \item $\cap$: intersection, \enquote{et} logique. \\
        \Ex $\{1; 2\} \cap \{2; 3\} = \{2\}$
    \item $\times$ et $E^{n}$: produit cartésien. \\
        \Ex \begin{itemize}
            \item Ensembles finis: $$
                \setf{1; 2} \times \setf{3, 4} 
              = \setf{(1; 3), (1; 4), (2; 3), (2; 4)}
            $$
            \item La droite réelle: $\mathbb{R}$,
                le plan réel: $\mathbb{R} \times \mathbb{R} = \mathbb{R}^2$
        \end{itemize}
\end{itemize}

\paragraph{Segments.} 

\begin{itemize}
    \item $\interval{a}{b}, \interval[open]{a}{b}$: 
        un intervalle ou \textit{segment} de $\mathbb{R}$.
    \item $\iinterval{a}{b}$: un intervalle d'entier. \\
        \Ex $\iinterval{2}{5} = \setf{2; 3; 4; 5}$
\end{itemize}


\subsection{Quantificateurs.}
\textsc{Attention}: en utilisant ses symboles, nous \textit{posons} une propriété.
Si une démonstration commence par un quantificateur, 
c'est que nous assumons cette propriété vraie.

\begin{itemize}
    \item $\forall$: \enquote{pour tout}, \enquote{quelque soit}. 
        C'est un quantificateur \enquote{puissant}:
        pour tout ses $x$ possible, les propriétés qui découlent 
        doit être vraies et prouvées. \\
        \Ex 
        \begin{equation}    
            \forall x \in \mathbb{R}, e^x > 0
        \end{equation}
        
    \item $\exists$: \enquote{il existe} au moins un. \\
        \Ex 
        \begin{equation}
            \forall y \in \mathbb{R}, \quad 
            \exists x \in \mathbb{R}, \quad
            y = x^2
        \end{equation}

    \item $\nexists$: \enquote{il n'existe aucun}. 
    \item $\exists!$: \enquote{il existe} un \textit{unique}.
\end{itemize}


\subsection{Divers.}

\paragraph{Nombres décimaux.}
\begin{itemize}
    \item $\floor{x} = \text{E}(x)$: partie entière, \enquote{sol} de $x$. \\
        \Ex $$
        \begin{aligned}
            \floor{x}  \leq \; &x < \floor{x} + 1 \\
            \floor{1.5}     &= 1 \\
            \floor{-4.1}    &= -5 \\
        \end{aligned}
        $$
    \item $\ceil{x}$: \enquote{plafond} de $x$. $\ceil{x} = \floor{x} + 1$
\end{itemize}

\paragraph{Analyse et suites.}
Pour un ensemble $E$.
Nous pouvons \textit{étendre} $E$: $$
    E \cup \setf{+\infty; -\infty}
$$.

\begin{itemize}
    \item $a \to b$ pour $a \in E, b \in E \cup \setf{+\infty; -\infty}$: 
        nous faisons tendre $a$ vers $b$. \\
        \textsc{Attention}: Tout de même, $a \neq b$.
    \item Pour $l \in E \cup \setf{+\infty; -\infty}$: 
        nous supposons l'existence d'une limite $l$ pour $f(a)$ quand $a$ tends vers $b$: 
        $$
            \lim_{a \to b}{f(a)} = l
        $$

    \item $(u_n)_{n > 0}, \forall n \in \mathbb{N}$: la suite $(u_n)$.
    \item Pour des ensembles $A$ et $B$, $f$ est un fonction qui à l'antécédent $x \in A$ 
        applique l'image $f(x) \in B$.
        $$
        \begin{aligned}
            f:\;&A &\to     \;\; &B \\
                &x &\mapsto \;\; &f(x)
        \end{aligned}  
        $$ 

    \end{itemize}

\end{document}
